% ==========================================================================
% GENERATED by scripts/split_paper.py. Do not edit this file.
%
% Source: imported_papers/grokking/iclr2027_conference_v2.tex
% Part 16: Adaptation to a new operation
%
% Edits belong upstream, in the published paper. To change what the
% chapter includes, change imports/grokking.toml and re-run
% the splitter.
% ==========================================================================
\section{Adaptation to a new operation}
\label{app:adapt}
This appendix measures adaptation speed, the training time to reach a new task from each grokked circuit. We grok on modular addition at $p=31$ at each setting of $\beta$, then fine-tune the whole network on a second operation with plain backpropagation, identically in every condition, counting epochs until test accuracy reaches $0.9$ (Table~\ref{tab:adapt}).

\begin{table}[!ht]
\caption{Full fine-tuning from a grokked circuit to a new operation, five or more seeds per condition, standard deviations in parentheses. Epochs count until test accuracy $0.9$ on the new operation. Free units are those whose outgoing norm falls below a tenth of the mean.}
\label{tab:adapt}
\begin{center}
\small
\begin{tabular}{lcccccc}
\toprule
Condition & Circuit PR & Free units & Source test & \multicolumn{2}{c}{Epochs to $0.9$} \\
 & & & & to sub & to mul \\
\midrule
$\beta=0$ (backprop) & $0.593\,(0.034)$ & $85$ & $0.995$ & $640\,(128)$  & $1800\,(141)$ \\
$\beta=0.25$         & $0.933\,(0.020)$ & $0$  & $0.989$ & $940\,(191)$  & $1940\,(233)$ \\
$\beta=0.5$          & $0.972\,(0.006)$ & $0$  & $0.990$ & $1200\,(89)$  & $1980\,(248)$ \\
$\beta=0.75$         & $0.983\,(0.003)$ & $0$  & $0.981$ & $1300\,(141)$ & $2020\,(133)$ \\
$\beta=1$            & $0.991\,(0.002)$ & $0$  & $0.972$ & $1540\,(180)$ & $1940\,(136)$ \\
From scratch         & ---              & ---  & ---     & $1410\,(137)$ & $1200\,(63)$ \\
\bottomrule
\end{tabular}
\end{center}
\end{table}

Adaptation to subtraction slows monotonically with the starting circuit's spread, at a rank correlation of $+1.00$ with participation, and the two ends do not overlap: every backpropagation seed lands between $400$ and $800$ epochs, every equalized seed between $1200$ and $1800$.

The free unit column separates two mechanisms. Backpropagation leaves about $85$ units carrying almost nothing, which a new task can claim instead of disturbing working structure. That count is already zero at $\beta=0.25$, yet adaptation keeps slowing all the way to $1540$ at $\beta=1$. Spare capacity therefore explains the jump from backpropagation to the distributed settings, and the continued rise among them means that the more evenly the work is spread, the more of it a new task has to disturb.

Multiplication is a different regime: every pretrained condition is slower than a fresh initialization, $1800$ to $2020$ epochs against $1200$, so what addition built gets in the way. The trade-off is therefore a claim about targets that benefit from the source task, which subtraction does and multiplication does not.

The direction is not special either. Subtraction to addition reproduces the effect, $740$ epochs at $\beta=0$ rising to $1600$ at $\beta=1$, again monotone at a rank correlation of $+1.00$.

Freezing the hidden circuit and retraining only the readout does not work at all: every condition stays at chance, at test $0.012$ or below. The units are not a reusable basis but combinations built for one operation, since units trained on $a+b$ compute quantities such as $\cos(\theta_a + \theta_b - \phi)$, and no linear combination of those gives the $\cos(\theta_a - \theta_b)$ that subtraction needs.
